\chapter{卫生宏}
\label{ch:19}

有些重复来自程序结构本身，而不仅是同一个函数被多次调用。宏在编译阶段组合语法，可以表达这种结构复用；代价是调用者需要更清楚地理解求值次数、名称来源和类型上下文。

本章先比较值参数与表达式片段，再分析结果类型推导、变参循环和卫生错误。例程均使用现有能力，尚未交付的声明生成与编译期字段遍历只作为边界说明，不作为可复制语法。

\section{值参数、引用与卫生}

先观察一个会改变状态的实参 next()。如果宏只是需要同一个结果两次，应该先取一次值再复用；如果宏接收的是待放置的表达式，两个放置位置则可能各执行一次。调用外观看起来相似，接口必须把这个差别说明白。

未标注参数采用单次求值的值参数规则：\$A(next()) 的一个实参一次，\$A(next(), next()) 的两个实参各一次。修改参数副本不改变调用方；只在保持行为时才能减少复制。显式 value 与 expr 的对比将在下一节用计数器展示。

显式 \&n 借用只读值，mut \&n 允许修改调用方；静态借用要求仍然生效。宏内部临时名保持卫生，不能从外部依赖某个展开时生成的临时名字。

\section{参数类别决定展开方式}

x:value 保留单次求值规则，x:expr 接收表达式语法并按展开位置执行：使用两次可能执行两次，未使用则不会执行。name:ident、T:type、body:stmt 和 body:block 分别要求相应语法类别。类别诊断发生在不匹配位置，不把任意语法都当作值参数。

多个 value 实参按调用实参顺序先求值；expr 片段随后随宏体中的放置位置执行。不要用 expr 传入有副作用的代码却假定它仍只执行一次。

\par
\medskip
\begin{lstlisting}[style=chali]
int64 count = 0;
int64 next() { count += 1; return count; }
$macro TwiceValue(x: value)
x + x
$end_macro
$macro TwiceExpr(x: expr)
x + x
$end_macro
void main() {
    int64 value = $TwiceValue(next());
    println("{} {}", value, count);
    count = 0;
    int64 expression = $TwiceExpr(next());
    println("{} {}", expression, count);
}
\end{lstlisting}
\par
\medskip

\par
\medskip
\begin{lstlisting}[style=out]
2 1
3 2
\end{lstlisting}
\par
\medskip

\subsection*{代码详解}

next 每调用一次便递增全局 count，再返回新的计数。TwiceValue 的 x:value 先取得一次 next() 的结果 1，再把同一个值相加得到 2，所以第一次输出 \texttt{2 1}，分别表示计算结果和真实调用次数。

把 count 清零后，TwiceExpr 将表达式语法放在两个使用位置，运行时 next() 执行两次，分别得到 1、2，相加为 3，因而输出 \texttt{3 2}。两种宏体看起来同样写 x+x，但参数类别改变求值方式；这个区别会直接影响输入读取、状态更新和资源取得，不能只看展开式的外观。

\subsection*{从求值次数选择参数类别}

先问宏要复用一个已经算出的值，还是把一段语法放到新的执行位置。读取输入、递增计数器或取得资源通常需要明确求值次数，优先考虑普通值参数。只有需要由展开后的控制结构决定执行位置时，才选择 expr，并向调用者说明它可能被跳过或重复执行。复制资源的成本与副作用也要按类型的复制规则判断，不能因为参数出现在宏里就忽略。

ident 传递名称，type 传递类型语法，stmt 与 block 传递语句或块；它们解决的是语法位置问题。类型参数 T 与形参 T:type 也不能仅凭字母相同就混为一谈。类别匹配之后，实际使用的名称、类型、借用与所有权仍须合法，传入一个标识符不会自动取得对该对象的写权限。

卫生意味着名字带有来源身份。模板自己声明的 result 可以与调用方的 result 同名而互不覆盖；模板中未声明的 temp 也不能靠调用方恰好存在一个 temp 来补全。这使宏可以在不同工程中复用，而不必要求每位使用者避开一份隐藏的临时名清单。需要影响调用方对象时，接口应明确表达这种关系。

\section{泛型参数与结果上下文}

宏可以声明类型参数。T:\#res 从使用结果的静态期望类型推导，U:b 从名为 b 的实参静态类型推导；它们不是在编译阶段运行实参表达式。显式类型参数优先，期望类型不足或约束冲突时需要补充类型。

例如 \$macro Example<T:\#res, U:b>(a, b) 同时使用结果与实参信息，推导出的 T、U 还必须满足宏体与调用位置的检查。

\par
\medskip
\begin{lstlisting}[style=chali]
$macro Example<T:#res, U:b>(a, b)
{
    T left = cast<T>(a);
    U right = b;
    yield left + cast<T>(right);
}
$end_macro
void main() {
    int32 result = $Example((int64)3, (int16)4);
    println("{}", result);
}
\end{lstlisting}
\par
\medskip

\par
\medskip
\begin{lstlisting}[style=out]
7
\end{lstlisting}
\par
\medskip

\subsection*{代码详解}

T:\#res 从 result 声明的 int32 取得期望结果类型，U:b 则从第二个实参 \texttt{(\allowbreak{}int16)\allowbreak{}4} 推导为 int16。第一个实参虽然是 int64，但宏体显式 cast<T>，把它转换到 left 的类型；right 先保留 U，再在相加位置转换成 T。

因此最终相加发生在符合结果上下文的类型中，值为 3+4=7。推导使用静态信息，不会为了知道 U 而预先执行 b。若去掉结果上下文或制造不相容约束，应提供显式类型或修正调用，不能把结果来源推导理解为运行期自动适配。

\section{标准宏与扩展边界}

使用现有标准宏构造常见结构时，仍需理解元素类型、求值次数和返回值。当前教程只使用已实现的类别和展开形式；声明宏尚未交付，不把声明生成、任意结果类别或未支持的重复匹配写成可用功能。

\par
\medskip
\begin{lstlisting}[style=chali]
$macro vec<T:first>[first,...rest]
{
    List<T> result;
    result.push(first);
    $for element in rest
    result.push(element);
    $end_for
    yield result;
}
$end_macro

void main()
{
    List<int64> result;
    result.push(99);
    List<int64> values = $vec[1, 2, 3];
    println("{} {} {}", values.len(), values.at(2), result.at(0));
}
\end{lstlisting}
\par
\medskip

\par
\medskip
\begin{lstlisting}[style=out]
3 3 99
\end{lstlisting}
\par
\medskip

\subsection*{代码详解}

vec 用方括号作为调用组，first 提供首个元素与 T 的推导来源，rest 收集后续实参。宏先创建自己的 List<T>，加入 first，再由编译期 \$for 为 rest 中每个元素生成 push 语句；这些 push 在生成的程序运行时执行，yield 交出列表。

调用 \texttt{\$vec[\allowbreak{}1,\allowbreak{} 2,\allowbreak{} 3]\allowbreak{}} 得到三项，at(2) 因为下标从 0 开始而取到 3。调用方另有 result，只含 99；卫生规则让宏内同名 result 不覆盖它，所以末项仍打印 99。不要把 \$for 与程序里的运行期 for 混为一谈，也不要把这套已有重复形式推广成尚未支持的任意模式匹配。

\section{诊断定位与检查}

宏报错时先区分定义位置与调用位置，检查类别、类型参数、展开作用域和借用。下面是需要被拒绝的卫生示例：宏内部临时变量不能作为调用方可见名字读取。

\par
\medskip
\begin{lstlisting}[style=chali]
$macro broken(value)
{
    temp.push(value);
    yield temp;
}
$end_macro

void main()
{
    List<int64> temp;
    List<int64> values = $broken(1);
}
\end{lstlisting}
\par
\medskip

\subsection*{代码详解}

broken 的模板直接使用 temp，却没有在模板自己的作用域中声明它，也没有通过形参接收这个来源。main 的 temp 虽然拼写相同，仍属于调用方；卫生名称解析不会把两者自动连接。

因此错误在编译检查阶段出现，没有正常运行输出。可以通过明确的参数接口表达需要的对象，或在宏内声明真正的局部列表，但不能靠调用方偷偷补一个同名变量修复模板。这里的 value 只提供元素值，没有提供列表来源。

\subsection*{展开与运行分别发生什么}

宏展开在编译阶段组合程序结构；示例中的 next 仍在程序运行时执行。编译器知道某个表达式的静态类型，不等于已经执行它。T:\#res 从赋值、返回等位置提供的期望类型取得信息，也不是运行程序以后观察结果再选择类型。缺少足够上下文时，应提供明确类型，而不是依赖偶然的字面量推导。

展开产物继续接受普通检查，之后进入四元式与 LLVM 后端。宏不能取消资源清理、绕过借用约束或给一个非法表达式增加执行资格。运行性能也取决于实际生成的操作及 LLVM 优化；少写模板代码不等于减少了运行工作量。定位错误时，可以先把调用的意图写成普通代码，确认类型和资源关系，再检查宏是否保留了这些关系。

本章的固定参数类别和已有变参循环是当前能力。显式结果类别、模块级声明生成、字段编译期遍历、member 宏和自定义派生扩展尚未交付，不能把设计中的形式复制为现行示例。运行期字段查询见下一章，它与编译时按字段生成 AST 的机制有不同的阶段和接口。

\subsection*{小结}

先用普通函数判断需求能否直接表达，再在确需语法组合时选择宏。阅读宏调用时，逐项标出哪些实参先得到值、哪些作为片段放置、名字属于调用方还是模板，以及结果类型来自哪里。

宏展开不执行任意用户运行期逻辑，也不使非法借用或资源使用合法。下一章的运行期反射解决按元数据访问对象的问题，不能与编译时生成语法结构混为同一机制。

\section{习题}

\begin{enumerate}

\item 比较 value 与 expr 中 next 的调用次数。

\item 修改未标注参数会改变调用方吗？

\item T:\#res 的类型来自运行期返回值吗？

\item 为什么不能读取宏内部 temp？

\end{enumerate}

\subsection*{参考解答与验收要点}

\begin{enumerate}

\item value 每个实参一次；expr 按实际放置和执行次数。

\item 不会，修改副本；需要修改调用方时显式可变引用。

\item 不是，来自静态期望结果类型。

\item 卫生隔离防止名称捕获，内部名不成为调用方声明。

\end{enumerate}

下一章：\hyperref[ch:20]{第 20 章}。

\section{番外：宏展开之后}

改写自《聊斋志异·画皮》。

太原王生，早行，遇一丽人，抱一卷宏模板，步甚艰。生悦而纳之，问其能何为。女曰：“妾可将三行变作百行，且不烦君亲写。”生大喜，遂与共处，连展开结果也不曾看。

一日，道士见生，曰：“子身上有两次求值之气。”生不信。归而窥之，见模板摊于屏上，next() 在两处各写一遍，却没有任何形参，拿笔者乃一副重复执行的骨架。

王生大惊，数其计数器，果然本欲加一，实已加二。女辩曰：“这两个 next() 本就是我在模板里写下的两次调用，怎能算作一个实参？”

生又在门外放一个 \texttt{temp}，指望宏内同名者来替自己持家。道士拂尘一挥：“卫生绑定，各有来处，莫拿隔壁同名之人顶替户主。”

生乃求救。道士令他改用一个表达式形参，编译器便先求值一次，再在模板中复用；需要修改外界之物则明写 mut \& 参数。展开之后，照旧过名称、类型与借用检查，不能披一张漂亮的宏皮便免检。

王生自此见到“少写九成代码”便先问：“生成的那九成，能给我看一眼么？”道士点头，说此人总算从审美转入了代码审查。

\subsection*{技术对照}

宏保留语法、来源与卫生绑定；表达式参数各求值一次并复用其值。普通参数修改的是副本，显式 mut \& 参数才写回来源。不能用调用者的同名变量补救模板中的未定义名称。
